%% GENERATED from PAPER_SUBMISSION.md (pandoc) -- hand-fixed conversion; see BUILD_NOTES.md
\section*{Appendix A --- Superseded earlier-stage supporting material (Beauty\_and\_\allowbreak Personal\_Care)}

\begin{quote}
\textbf{Why this is an appendix:} the material below is an earlier-stage Beauty\_and\_\allowbreak Personal\_Care reproducibility study (the 20-variant cross-pipeline scan, its 2-seed signal, and the retracted LIGER-gap audit). It predates and does not belong to the headline spine (causal FIR filter + dataset-conditional tail pattern + competitive-overall Video\_Games result, \S\ref{sec:1}--\S\ref{sec:7}). It is retained verbatim for the reproducibility and negative-result record only, and is referenced from \S\ref{sec:6.1}--\S\ref{sec:6.2} as supporting context. The live Beauty result that enters the tail pattern (\S\ref{sec:5.3}) is the text\ensuremath{-}ID tail contrast, not these single-seed sampled-softmax scans.
\end{quote}


\subsection*{A.1 Beauty\_and\_\allowbreak Personal\_Care --- 20-variant cross-pipeline scan (superseded supporting material)}

\begin{quote}
\textbf{Note:} the following Beauty\_and\_PC 20-variant scan and the LIGER-gap audit are \textbf{earlier-stage supporting material}, retained for the reproducibility and negative-result record. They predate the HSTU + causal-filter + tail-pattern spine and are \emph{not} the paper's headline; the live Beauty result used in the tail pattern (\S\ref{sec:5.3}) is the text\ensuremath{-}ID tail contrast, not these single-seed sampled-softmax scans.
\end{quote}

The scan varies five axes: \textbf{text encoder} (MiniLM titles/rich-text, BLaIR titles/rich-text), \textbf{projection \ensuremath{\Phi}} (single Linear vs the 2-layer MLP adaptor of \citealp{hou2024blair}), \textbf{learned item embedding} (kept vs removed = faithful SASRecText), \textbf{loss} (sampled softmax K=1024, in-batch negatives, chunked-full-softmax, hybrid), and \textbf{augmentation} (random subsequence sampling, factor 1/2/3). Training used a 15--20-epoch budget, batch size 256, Adam (lr 1e-3, weight decay 1e-5, gradient clip 5.0), a single seed for the initial scan, and two seeds (42, 43) for the final reported variants. Table A1 summarizes all 20 variants. \textbf{Numbers are single-seed unless otherwise noted.}

%% GENERATED by _bestrec_run/emit_latex_tables.py -- DO NOT EDIT BY HAND
% table key: tableA1
% provenance: md lines 777-800; family=md-only (pandoc-converted; no JSON family); check={"checked": 0, "note": "md-only table; no JSON family"}
\par\medskip\begingroup\footnotesize\setlength{\tabcolsep}{4pt}\renewcommand{\arraystretch}{1.12}
\noindent \textbf{Table A1: Beauty\_and\_\allowbreak PC 5-core cross-pipeline scan (full-catalog eval, NDCG@10) --- 22 numbered rows: 20 distinct variants plus two seed replications (row 20 replicates variant 19; row 22 replicates baseline variant 7)}\par\smallskip
\noindent
\begin{tabularx}{\linewidth}{cYrr}
\toprule
\# & Variant & NDCG@10 & \ensuremath{\Delta} vs base \\
\midrule
1 & Popularity & (not run) & --- \\
2 & SASRec sampled-1024 + MiniLM titles & 0.0180 & --- \\
3 & SASRec sampled-1024 + MiniLM titles, d=128 & 0.0176 & -2.2\% \\
4 & SASRec in-batch negs + augment-2 & 0.0036 & -80\% \\
5 & TIGER minimal (our re-impl) & 0.0080 & -56\% \\
6 & SASRec hybrid (in-batch + 8K random negs + aug-3) & 0.0106 & -41\% \\
7 & \textbf{SASRec chunked-full + MiniLM titles (improved baseline)} & \textbf{0.0191} & \textbf{+6.1\%} \\
8 & SASRec chunked-full + MiniLM titles, d=128 & 0.0186 & -2.6\% (vs 7) \\
9 & SASRec chunked-full + MiniLM titles, 25 epochs & 0.0189 & -1.0\% (vs 7) \\
10 & BERT4Rec & 0.0160 & -16\% (vs 7) \\
11 & SBERT-only (no learned item\_emb) + MiniLM & discontinued & --- \\
12 & Faithful SASRecText (no item\_emb + MLP + dropout 0.5 + BLaIR) & discontinued @ epoch 2 & --- \\
13 & Hybrid (item\_emb + MLP + dropout 0.5 + BLaIR) & discontinued @ epoch 5 & --- \\
14 & BLaIR-titles + Linear projection (just encoder swap) & discontinued @ epoch 2 & --- \\
15 & Fusion (BLaIR + rich-text MiniLM concat) + Linear & discontinued @ epoch 1 & --- \\
16 & \textbf{BLaIR-titles + MLP adaptor (\citet{hou2024blair}) + dropout 0.2} & \textbf{0.01928} & \textbf{+0.9\% (vs 7)} \\
17 & + augment-2 & 0.0036 val & discontinued; aug hurts \\
18 & + d=128 width & discontinued at eval (slow) & --- \\
19 & \textbf{BLaIR-rich-text + MLP adaptor + dropout 0.2 (seed 42)} & \textbf{0.01936} & \textbf{+1.3\% (vs 7)} \\
20 & seed 43 of variant 19 & \textbf{0.01955} & \textbf{+2.3\% (vs 7)} \\
21 & \textbf{MiniLM + MLP adaptor (clean ablation)} & \textbf{0.01889} & \textbf{-1.2\% (vs 7)} \\
22 & Baseline seed 42 (reproduction) & 0.0190 & 0.0\% (vs 7) \\
\bottomrule
\end{tabularx}
\endgroup\par\medskip


\textbf{2-seed mean of variant 19+20 (BLaIR-rich-text + MLP)}: NDCG@10 = 0.01946, range 0.0002.

\textbf{Clean-ablation finding (variant 21, added after audit)}: When we hold the encoder constant at MiniLM and ONLY change projection layer from single-Linear to the SASRecText 2-layer MLP adaptor, NDCG@10 drops slightly from \textasciitilde0.0190 to 0.01889. \textbf{The MLP adaptor by itself does NOT help on this protocol.} The +1.3-2.3\% gains we observe with variants 16/19/20 are attributable to the change in text encoder (MiniLM \ensuremath{\rightarrow} BLaIR) and the change in text content (titles \ensuremath{\rightarrow} rich text), not to the projection layer.

Re-analysis of the 20-variant scan:

\begin{itemize}
\tightlist
\item
  \textbf{MLP adaptor alone}: no benefit (variant 21 vs 7 = -1.2\%)
\item
  \textbf{MiniLM \ensuremath{\rightarrow} BLaIR encoder}: small benefit (variant 16 vs 7 = +0.9\% NDCG@10, single-seed)
\item
  \textbf{Titles \ensuremath{\rightarrow} rich text (BLaIR)}: small additional benefit (variant 19 vs 16 = +0.4\%, single-seed)
\item
  \textbf{Removing learned item embedding} (faithful SASRecText): hurts substantially
\item
  \textbf{Heavy dropout (0.5)} as in the published SASRecText: hurts on our 5-core data (median 5 interactions/user)
\item
  \textbf{Subsequence augmentation}: hurts (\textasciitilde-5\% NDCG)
\item
  \textbf{Width (d=128)}: no help
\end{itemize}


\subsection*{A.2 Beauty\_and\_PC multi-seed signal (superseded supporting material)}

\begin{quote}
\textbf{Note:} this is appendix supporting material for the Beauty 20-variant scan (\S{}A.1), not the headline tail-pattern result (\S\ref{sec:5.3}).
\end{quote}

We have multi-seed numbers for two configurations:

\begin{itemize}
\tightlist
\item
  Best variant (BLaIR + rich-text + MLP): seed 42 \ensuremath{\rightarrow} 0.01936, seed 43 \ensuremath{\rightarrow} 0.01955; \textbf{mean 0.01946 (n=2)}.
\item
  Baseline (chunked-full, no MLP, MiniLM titles): original \ensuremath{\rightarrow} 0.01911, seed 42 \ensuremath{\rightarrow} 0.0190; \textbf{mean 0.01906 (n=2)}.
\end{itemize}

Apparent gain: 0.01946 - 0.01906 = \textbf{+0.0004 (+2.1\% relative)} with n=2 vs n=2. Inter-seed variance for our baseline (0.0001 range) is roughly half the size of the gain, suggesting the directional signal is real but small. A formal significance test would require larger n; our paper-stated +2.1\% improvement is honest about this.


\subsection*{A.3 Audit of ``gaps to LIGER'' claims (corrected from earlier versions; superseded supporting material)}

\begin{quote}
\textbf{Note:} this comparator-audit retraction is retained as appendix material; the body's \S\ref{sec:5.4} is the titration result.
\end{quote}

Earlier versions of this work cited an apparent 57\% gap to LIGER on Beauty\_and\_\allowbreak Personal\_Care. \textbf{After a comparator audit, this framing is retracted.} The facts:

\begin{enumerate}
\def\labelenumi{\arabic{enumi}.}
\item
  \textbf{LIGER reports on Amazon 2014, not AR2023.} The LIGER paper \citep{yang2024liger} evaluates on Amazon Beauty (2014, \citealp{he2016ups}), reporting NDCG@10 = 0.04020 \ensuremath{\pm} 0.00044 (K=20) and 0.04738 \ensuremath{\pm} 0.00151 (K=N) on their 5-core preprocessing. The 2014 Amazon Beauty subset has tens of thousands of items, while our AR2023 Beauty\_and\_\allowbreak Personal\_Care 5-core has 207,649 items. \textbf{These are different datasets.} The ``57\% gap'' framing in earlier versions compared our AR2023 number to LIGER's 2014 number, which is invalid.
\item
  \textbf{There is no published LIGER number on AR2023 Beauty\_and\_\allowbreak Personal\_Care 5-core.} LIGER does not evaluate on AR2023 at all in the published paper.
\item
  \textbf{TIGER reports on Amazon 2014 only.} The TIGER paper \citep{rajput2023tiger} reports NDCG@10 = 0.0384 on Amazon Beauty 2014 with their 5-core, also not on AR2023.
\item
  \textbf{BLaIR \citep{hou2024blair} introduced AR2023 --- historical note, superseded: the AR2023 5-core literature has since grown (HSTU-BLaIR and 2026 preprints; \S\ref{sec:5.1}).} Their Beauty subset is ``All\_Beauty'' (not ``Beauty\_and\_\allowbreak Personal\_Care'') and their preprocessing is by-timestamp 8:1:1 with no k-core filter. They report NDCG@10 on All\_Beauty in the range 0.0177-0.0241 across various LLM encoders (Table 8). With a UniSRec downstream model, their best All\_Beauty number is 0.0241 (gemini-embedding). They do not report on Beauty\_and\_\allowbreak Personal\_Care.
\end{enumerate}

We therefore cannot make a defensible gap claim to TIGER/LIGER/BLaIR on AR2023 Beauty\_and\_\allowbreak Personal\_Care 5-core, because \textbf{we did not find a comparable published number} (a search statement as of this writing, not a priority claim --- AR2023 preprints are appearing rapidly; \S\ref{sec:5.1}). We provide our 0.01946 (2-seed mean) as a reference baseline rather than as a SOTA-improvement claim.

A future apples-to-apples comparison would require either:

\begin{itemize}
\tightlist
\item
  Running TIGER/LIGER on our 5-core preprocessing (multi-week effort for faithful reproduction), or
\item
  Adopting Hou et al.'s by-timestamp 8:1:1 preprocessing and re-running our pipeline.
\end{itemize}

